%% Согласно ГОСТ Р 7.0.11-2011:
%% 5.3.3 В заключении диссертации излагают итоги выполненного исследования, рекомендации, перспективы дальнейшей разработки темы.
%% 9.2.3 В заключении автореферата диссертации излагают итоги данного исследования, рекомендации и перспективы дальнейшей разработки темы.
\begin{enumerate}
  \item При создании гибридного источника излучения на длине волны 3 мкм с волоконной накачкой в маломодовом ОВ доставки впервые обнаружено взаимное влияние ОМЧВС и ММЧВС, имеющих общую антистоксовую компоненту на длине волны 1003~нм. Показано, что синхронизация импульсов накачки (1030~нм) и затравки (1562~нм) увеличивает спектральную мощность компонент ММЧВС на 20 дБ, что впоследствии снижает эффективность ГРЧ в н-о кристалле. Предложен и реализован метод подавления взаимного влияния ЧВС с помощью модового фильтра на входе ОВ, позволяющий увеличить допустимую длину ОВ доставки с 2.5 м до 3.0 м без существенной потери полезной мощности накачки.
  \item Впервые экспериментально исследована эволюция КФС при фотоиндуцированной ГВГ в германосиликатном ОВ в режиме самозатравки. Показано, что в процессе оптического полинга ширина КФС и их сдвиг монотонно уменьшаются до стационарных значений, а форма кривых сохраняет выраженную асимметрию. Проведён соответствующий численный расчёт с использованием модели~\cite{antonyuk2001self} с добавлением учёта керровского эффекта. Обнаружено, что на поздних стадиях эффективность ГВГ остаётся практически постоянной вследствие динамического равновесия между уменьшением спектральной расстройки и ростом пиковой эффективности преобразования. Полученные результаты создают основу для разработки физической модели фотоиндуцированной ГВГ в ОВ, которая позволит прогнозировать уровень паразитной ГВГ в ОВ доставки мощных лазеров, в которых данный эффект может быть нежелательным, поскольку он может приводить к снижению порога оптического разрушения н-о кристаллов при ГРЧ в гибридных лазерных системах.
  \item Предложены поправочные выражения для трёх методик обработки температурных кинетик при разогреве образца лазерным излучением в методе ЛК, позволяющие компенсировать нестабильность температуры окружающей среды и коэффициента теплоотдачи. Экспериментально показано, что применение поправок снижает разброс измеренных значений коэффициента оптического поглощения более чем в два раза (с 23\% до 11\%).
\end{enumerate}
